\chap{Abstract} 


%%%% your abstract goes here (word limit: 350)
 A large amount of scientific research is funded by the public, and therefore the public has a right to access scientific discovery.  Approximately 40\% of basic research in the United States is government-funded~\cite{Pece_Anderson_2024}. To promote equitable access to publicly funded research, the U.S. National Institutes of Health now requires all NIH-funded papers to be open access.\footnote{\url{https://grants.nih.gov/policy-and-compliance/policy-topics/public-access/nih-public-access-policy-overview}} However, open access alone does not guarantee comprehension. Effective scientific summarization must consider multiple factors, including both form and audience For example, a summary for a trained biologist might differ from one for a high school student taking biology. Even on the user level, an individual's informational needs may change depending on the circumstance. The same trained biologist may require a longer, more technical summary when conducting a replication study, but may prefer a shorter, more concise summary when conducting a general literature survey. Current summarization methods often make assumptions about the intent and desires of the end user in order to constrain the requirements for building the system. However, these assumptions also constrain the potential utility the system poses to the user. Summarization models must account for user needs and circumstance, such as differing document length, background, or technical level. \textbf{My research reduces barriers to accessing scientific knowledge by developing methods to generate and evaluate summarization systems that focus on the informational and practical needs of the end user.}



%% list of keywords seperated by comma
\keywords{summarization, natural language processing, human-computer interaction, evaluation, natural language generation, language models, artificial intelligence}


%%%%  committee members (add it right after the abstract w/o page break)
\begin{singlespace}

    %% if you have co-advisor, add here w/ \vspace{0.1in} as shown below
    %% alternatively you can use minipage environment to put side-by-side
    \section*{Primary reader and thesis advisor}
    
    Dr. Mark Dredze \\
    Professor\\
    Department of Computer Science\\
    Johns Hopkins University, Baltimore MD 


    \section*{Secondary readers}
    
    Dr. Daniel Khashabi\\
    Assistant Professor\\
    Department of Computer Science \\
    Johns Hopkins University, Baltimore, MD 
    
    \vspace{0.1in}
    
    Dr. Sharon Levy \\
    Assistant Professor\\
    Department of Computer Science \\
    Rutgers University, New Brunswick, NJ

    %% you can add more readers if you have them on your committee 
    %% use \vspace{0.1in} in between members for clarity
    %% you can also place committee members side-by-side using `minipage`


\end{singlespace}